\section{CLIP}
\textbf{CLIP} (Contrastive Language-Image Pretraining) is a state-of-the-art \textbf{Vision-Language Model (VLM)} developed by OpenAI. 
It leverages natural language supervision to learn a wide range of visual concepts by simply predicting the correct image-caption pairings during pre-training.

The architecture consists of two main components:
\begin{itemize}
    \item An \textbf{image encoder} (typically a CNN or Vision Transformer).
    \item A \textbf{text encoder} (typically a standard Transformer).
\end{itemize}

CLIP's key innovation is training both encoders simultaneously using a \textbf{contrastive learning approach}. 
This mechanism maps images and text into a single, shared embedding space, forcing semantically similar pairs close together while pushing dissimilar pairs apart. 
This unified representation is what allows the model to easily perform versatile downstream tasks, such as zero-shot learning, image classification, and object detection.

\paragraph{Before CLIP}
Prior to CLIP, traditional computer vision relied on heavily curated datasets of images and fixed, categorical labels. This supervised approach had significant limitations:

\begin{itemize}
    \item \textbf{Data collection bottleneck:} Acquiring large volumes of high-quality, human-annotated data is expensive and time-consuming.
    \item \textbf{Poor generalization:} Models were confined to their specific training tasks and struggled to adapt to new domains.
    \item \textbf{Inflexible classes:} Introducing even a single new class required inefficiently retraining the model from scratch.
\end{itemize}

By replacing fixed labels with open-ended textual descriptions, CLIP overcomes these hurdles, enabling zero-shot classification of novel visual concepts without any task-specific retraining.

\subsection{Ideas that led to CLIP}
Rather than introducing an entirely new paradigm, CLIP successfully combined and scaled existing machine learning concepts:
\begin{itemize}
    \item \textbf{Webly supervised learning}: Exploiting the vast, readily available pools of data on the internet to train models without manual annotation.
    \item \textbf{Weak supervision from multimodal data}: Leveraging naturally occurring pairings of visual and textual information to learn a joint embedding space.
    \item \textbf{Scaling laws}: Natural Language Processing (NLP) had already demonstrated massive performance leaps by scaling up models and datasets (e.g., GPT-3). CLIP proved this principle translates effectively to computer vision.
\end{itemize}

\subsection{Dataset: WebImageText (WIT)}
Previous vision-language research relied on relatively small datasets like MS-COCO, which contains roughly 100,000 images. To test whether their approach could truly scale, OpenAI constructed a massive new dataset called \textit{WebImageText} (WIT).

WIT consists of 400 million image-text pairs scraped from the internet. To ensure diversity and broad conceptual coverage, the collection process utilized 500,000 distinct search queries encompassing common terms, word combinations, and random concepts. To maintain an approximate class balance, the dataset was capped at 20,000 image-text pairs per query.

While WIT is orders of magnitude larger than prior datasets, it is also inherently noisy. Because the web data was minimally curated, the quality of the image-text pairings is not strictly guaranteed, meaning the resulting model inevitably inherits the biases and limitations of the internet.

\subsection{Contrastive Pretraining}
The core objective of CLIP's training phase is to \textbf{maximize} the cosine similarity between the embeddings of the $N$ matching image-text pairs in a batch, while \textbf{minimizing} the similarity of the $N^2 - N$ incorrect pairings.

This is achieved using a symmetric, softmax-based contrastive loss function:
\begin{equation}
    \loss = -\frac{1}{2N} \sum_{i=1}^{N} \left(
    \underbrace{\log \frac{\exp(\simop(\vct{v}_i, \vct{t}_i) / \tau)}{\sum_{j=1}^{N} \exp(\simop(\vct{v}_i, \vct{t}_j) / \tau)}}_{\textit{Image} \to \textit{Text}}
    +
    \underbrace{\log \frac{\exp(\simop(\vct{t}_i, \vct{v}_i) / \tau)}{\sum_{j=1}^{N} \exp(\simop(\vct{t}_i, \vct{v}_j) / \tau)}}_{\textit{Text} \to \textit{Image}}\right),
\end{equation}
where $\vct{v}_i$ and $\vct{t}_i$ represent the embeddings of the $i$-th image and text pair, and $\simop$ is the cosine similarity function.

The two terms of the sum play a symmetric but distinct role:
\begin{itemize}
    \item \textbf{Image} $\to$ \textbf{Text}: for each image embedding $\vct{v}_i$, this term is a softmax over all $N$ text embeddings in the batch. It pulls $\vct{v}_i$ closer to its matching text embedding $\vct{t}_i$, while pushing it away from all the other $N-1$ text embeddings.
    \item \textbf{Text} $\to$ \textbf{Image}: symmetrically, for each text embedding $\vct{t}_i$, this term is a softmax over all $N$ image embeddings, pulling $\vct{t}_i$ closer to its matching image embedding $\vct{v}_i$, while pushing it away from all the other $N-1$ image embeddings.
\end{itemize}
Averaging both directions ensures the embedding space is aligned consistently, regardless of whether an image is used to query a text or vice versa.

\termdef{Temperature}{The hyperparameter $\tau$ modulates the sharpness of the probability distribution. A lower temperature sharpens the distribution, which boosts confidence but increases the risk of overfitting. Conversely, a higher temperature softens the distribution, regularizing the model at the potential cost of peak performance.}

\termdef{Training Phase}{Due to the immense scale of the WebImageText dataset, both the image and text encoders are trained entirely from scratch without relying on pre-trained weights. Only basic data augmentation techniques are applied to the visual inputs during this process.}

\paragraph{Limitations of the training approach}
\label{CLIP:training-limitations}
This training approach has some limitations regarding the computational resources required. 
In particular, since we need to compute the cosine similarity between all pairs of images and text in a batch ($N^2$ pairs),
the size of the batch is limited by the available memory.

\subsection{Inference}
During inference, CLIP can be used for various tasks such as image classification, object detection, and zero-shot learning.
\begin{itemize}
    \item \textbf{Image classification}: Given an image, CLIP can be used to classify it by comparing its embedding to the embeddings of a set of candidate labels (text descriptions);
    \item \textbf{Text to Image retrieval}: Given a text query, CLIP can be used to retrieve relevant images by comparing the text embedding to the embeddings of a set of candidate images.
    First, the text query is encoded using the text encoder to obtain its embedding. Then, the cosine similarity between the text embedding and the embeddings of a set of candidate images is computed. 
    The images with the highest similarity scores are returned as the most relevant results.
    \item \textbf{Image to Image retrieval}: Given an image query, CLIP can be used to retrieve relevant images by comparing the image embedding to the embeddings of a set of candidate images.
    Similarly to text to image retrieval, the image query is encoded using the image encoder to obtain its embedding. Then, the cosine similarity between the image embedding and the embeddings of a set of candidate images is computed.
    The images with the highest similarity scores are returned as the most relevant results.
    \item \textbf{Zero-shot learning}: CLIP enables zero-shot learning, meaning it can recognize and classify images based on textual descriptions without needing to be retrained on specific classes.
\end{itemize}

\subsection{Fine Tuning}
Once the power of CLIP was demonstrated on zero-shot learning tasks, the next step was to explore how to further improve its performance on specific tasks.

A first approach was to fine-tune either the image encoder, the text encoder, or both on a specific dataset.
This was found to be ineffective, as the model had already learned a powerful shared embedding space, and fine-tuning on a specific dataset often led 
to overfitting and a decrease in performance on other tasks.

Other approaches based on prompt engineering were found to be more effective in improving the performance of CLIP on specific tasks, 
without the need for fine-tuning the model itself.
Here, we will discuss some of the most effective techniques for improving CLIP's performance through prompt engineering.

\subsubsection{Basic Prompt Engineering}
Prompt engineering is the process of designing and optimizing the input prompts to improve a model's performance on specific tasks. 
Since CLIP learns a shared embedding space, the phrasing of the text prompts significantly impacts the results.

For example, since images are rarely described with a single word in the wild, using a template like \texttt{"a photo of a [class]"} 
often outperforms using the raw class name alone. On ImageNet, this simple modification yielded a gain of 1.3\% in accuracy.

\subsubsection{Prompt Ensembling}
Further improvements can be achieved by using multiple prompts for each class and averaging their embeddings. 
For the class "dog," one might ensemble prompts like \textit{"a photo of a dog,"} \textit{"a picture of a dog,"} and \textit{"a photo of a cute dog."} On ImageNet, 
this approach provided an additional 3.5\% gain in accuracy compared to a single prompt.

\subsubsection{Advanced Approaches}
Prompt engineering can still be difficult and time-consuming. Many times a slight change in the prompt can lead to a significant change in performance,
with no clear way to predict which prompts will work best for a given task.

To address this, more advanced approaches have been proposed.

\paragraph{CoOp}
CoOp (Context Optimization) is a method that learns continuous \textbf{prompt vectors} instead of using discrete text prompts.
These prompt vectors are then fed into the text encoder to generate the class embeddings.
Different choices can be made, such as the number of prompt vectors, their dimensionality, where they are inserted in the text encoder (e.g., before the class name, after the class name, or both), 
and whether they are shared across classes or specific to each class.

CoOp optimizes the prompt vectors using a small amount of labeled data from the target task, allowing it to adapt the prompts to better suit the specific task at hand.
This approach has been shown to significantly improve the performance of CLIP on various tasks, especially when the amount of labeled data is limited.

Even though the results are quite impressive, they are still not humanly interpretable, as the learned prompt vectors may not correspond to any meaningful words or phrases.
Also, the approach was shown to have poor generalization performance to new classes, as the learned prompt vectors may be overfitted to the specific classes seen during training.

\paragraph{CoCoOp}
CoCoOp (Contextualized Context Optimization) is an extension of CoOp that addresses the issue of generalization to new classes. 

CoCoOp learns a \textbf{meta-network} that generates some meta-prompts based on the output of the image encoder. 
These meta-tokens are then summed with the learned prompt vectors to generate the final prompts that are fed into the text encoder.
This allows CoCoOp to adapt the prompts based on the specific image being processed, improving generalization to new classes and reducing overfitting to the classes seen during training.

The problem with this approach is that it is more computationally expensive than CoOp, as it requires an additional forward pass through the meta-network for each image,
during the \textbf{training phase}.

\paragraph{Visual Prompt Tuning}
All the approaches described so far have focused on optimizing the text prompts, while keeping the image encoder fixed.
Visual Prompt Tuning is a method that optimizes the input to the image encoder instead of the text encoder. 
Similarly to CoOp, it learns a set of continuous prompt vectors that are instead added to the input image before being fed into the image encoder.

\paragraph{Multimodal prompt Learning}
Sometimes, optimizing only the text prompts or only the image prompts leads to suboptimal performance, as ther's no general rule to determine which one is more effective for a given task.
Multimodal Prompt Learning is a method that optimizes both the text and image prompts simultaneously, allowing for a more holistic approach to prompt engineering.

Using a small neatwork (for example, a transformer), the method learns to generate both text and image prompts based on the input image and the target task.
The prompts are then fed into the respective encoders to generate the final embeddings, which are used for classification or other downstream tasks.

\subsection{Beyond CLIP - SigLIP}
To resolve the limitations of CLIP, presented in \autoref{CLIP:training-limitations}, a new method called SigLIP (Sigmoid Loss for Language-Image Pretraining) was proposed.
Instead of using a softmax-based contrastive loss, where each pair of image and text is compared to all other pairs in the batch, 
SigLIP uses a sigmoid-based loss that compares each pair independently, treating the problem as a binary classification.

For each pair of image and text, the model predicts a score that indicates whether the pair is a match or not (e.g., whether the image and text are semantically related).
Because each pair is treated independently, the batch size can be much larger without running into memory issues, allowing for more efficient training on large datasets.

